\lesson{2}{Nov 15 2021 Mon (12:15:15)}{Polynomial Transformations}{Unit 3}

\begin{definition}[Theorems]
    A theorem is a proven statement based on experimentation. I will go over
    these three \textbf{Theorems}:
    
    \begin{itemize}
        \item Fundamental Theorem of Algebra
        \item Factor Theorem
        \item Remainder Theorem
    \end{itemize}
\end{definition}

\begin{theorem}[Fundamental Theorem of Algebra]
    The Fundamental Theorem of Algebra generally states that the degree of a polynomial is equivalent to the number of zeros (both real and complex) of a function.
    
    By the \textbf{Fundamental Theorem of Algebra}, the polynomial function $f(x) = x^2 - 3x - 28$ has two zeros since the degree of the function is two. To determine these zeros, replace the function notation of $f(x)$ with $0$ and solve by factoring.
    
    To factor, let's use the \textbf{Quadratic Equation}:
    
    \begin{align}
        x    &= \frac{-b \pm \sqrt{b^2 - 4ac}}{2a} \\
        f(x) &= x^2 - 3x - 28 \\
        a    &= 1, b = -3, c = -28 \\
        x    &= \frac{-(-3) \pm \sqrt{(-3)^2 - 4 \times 1(-28)}}{2 \times 1} \\
        x    &= \frac{3 + 11}{2} \rm{ AND } x = \frac{3 - 11}{2} \\
        x    &= 7 \rm{ AND } -4 \rightarrow (x + 4)(x - 7)
    \end{align}
    
    The zero product property tells us that for these factors to result in a product of $0$, one or both of them must equal $0$.
    
    \begin{minipage}[c]{0.5\textwidth}
        \centering
        \begin{tabular}{cccccc}
            & 0 & = & x & + & 4 \\
            - & 4 & = &   & - & 4 \\
            \hline
            & x & = & & - & 4 \\
        \end{tabular}
    \end{minipage}
    
    \hfill
    \begin{minipage}[c]{0.5\textwidth}
        \centering
        \begin{tabular}{cccccc}
            & 0 & = & x & - & 7 \\
            + & 7 & = &   & + & 7 \\
            \hline
            & x & = & & + & 7 \\
        \end{tabular}
    \end{minipage}
    
    The \textbf{Zeros} of $f(x) = x^2 - 3x - 28$ are: $x_1 = -4, x_2 = 7$.
\end{theorem}

\begin{marginfigure}
    \centering
    \incfig{fundamental-theorem-of-algebra}
    \sidecaption{$f(x) = x^2 - 3x - 28$ Graphed.}
    \label{fig:fundamental-theorem-of-algebra}
\end{marginfigure}

\begin{theorem}[Factor Theorem]
    The \textbf{Factor Theorem} states that a first degree binomial is a factor of a polynomial function if the remainder, when the polynomial is divided by the binomial, is $0$.

    To determine whether $x - 5$ is a factor of the function $f(x) = -4x^3 + 21x^2 - 25$, set up a division problem whereby $-4x^3 + 21x^2 - 25$ is divided by $x - 5$.
    
    \[ \polylongdiv{-4x^3 + 21x^2 - 25}{x - 5} \]

    When the function $f(x) = -4x^3 + 21x^2 - 25$ is divided by the binomial $x - 5$, the remainder is $0$. So, $x - 5$ is a factor of the function $f(x) = -4x^3 + 21x^2 - 25$.
\end{theorem}

\begin{theorem}[Remainder Theorem]
    The \textbf{Remainder Theorem} states that when the opposite of the constant from the binomial divisor is substituted into a function for $x$, the result is the remainder.

    When the polynomial function $f(x) = x^4 + 11x^3 + 26x^2 + 15x - 17$ is divided by $x + 8$ using division, the remainder is the last integer on the bottom row.

    \[ \polylongdiv{x^4 + 11x^3 + 26x^2 + 15x - 17}{x + 8} \]
    
    When the opposite constant in the divisor is substituted into the function, the result will be the same as the remainder in the division process.
    
    \begin{align}
        f(x)  &= x^4 + 11x^3 + 26x^2 + 15x - 17 \\
        f(-8) &= (-8)^4 + 11(-8)^3 + 26(-8)^2 + 15(-8) - 17 \\
              &= 4096 + 11(-512) + 26(64) + 15(-8) - 17 \\
              &= 4096 - 5632 + 1664 - 120 - 17 \\
              &= -9 \leftarrow \rm{Remainder}
    \end{align}
\end{theorem}

\subsubsection*{Using these Theorems}

Given a polynomial function $f(x)$ and a number $a$, if $(x - a)$ is a factor of $f(x)$, then $a$ is a $0$ of the polynomial.

The binomial $(x - a)$ can be proved as a factor of $f(x)$ by:

\begin{itemize}
    \item Using \textbf{Long Division} with $(x - a)$ as the divisor.
    \item Using factoring methods when appropriate (grouping, completing the square, \ldots).
\end{itemize}

If a is a $0$ of the polynomial function $f(x)$, then:

\begin{itemize}
    \item The graph of $f(x)$ crosses the $x$ axis at $(a,0)$.
    \item Substituting $a$ into $(x - a)$ will equal $0$
    \item Substituting $a$ into $f(x)$ will result in $f(a) = 0$.
\end{itemize}

\newpage
